\documentclass[11pt,a4paper]{article}
\usepackage[margin=2cm]{geometry}
\usepackage{amsmath}
\usepackage{xcolor}
\usepackage[colorlinks=true,urlcolor=blue!50!black]{hyperref}
\usepackage{microtype}
\pagestyle{empty}
\setlength{\parskip}{4pt}
\setlength{\parindent}{0pt}
\begin{document}
{\Large\bfseries Twisted Kerr-mass spacetimes: evidence that orbits along the axis lose Kerr's hidden symmetry}\\[2pt]
{\small Plain-language summary of a draft preprint by Sumit Pal (independent researcher). AI-assisted; not yet
reviewed by an outside expert. 28 November 2026.}

\textbf{The question.} Orbits around a spinning (Kerr) black hole are unusually orderly. Besides energy and angular
momentum, they conserve a hidden quantity, Carter's constant, so they never become chaotic. Is that orderliness
fragile? We take a family of spacetimes that look exactly like Kerr in their mass distribution along the spin axis,
but whose ``twist'' (how the space is dragged around) is changed by an amount $k$. $k=0$ is Kerr; physically sensible
members run from $k=0$ to $k=30/7\approx4.29$.

\textbf{What we did.} We followed one orbit, a particle falling up and down the spin axis, and asked whether nearby
orbits keep enough conserved quantities to stay orderly. A theorem from differential Galois theory (Morales--Ramis)
turns this into a test on a single matrix group. If the group is as large as possible, there cannot be enough
well-behaved conserved quantities.

\textbf{What we found.}
\begin{itemize}\itemsep2pt
\item \textbf{Numerically,} for every twist in the physical range ($0<k\le30/7$), at four spin values and one
energy: the group is as large as possible. So, at this numerical grade, Carter's constant does not survive any amount of this twist.
\item \textbf{Proved,} with explicit, computer-checked margins, very close to Kerr ($0<k\le$ a few
thousandths), given three stated assumptions (one of them our own unrefereed proof).
\item \textbf{Certified by computer-assisted interval arithmetic} (every rounding error bounded) over most of the
range at two spins: $0<k\le3$ at spin $0.6$; $0.003\le k\le3$ at spin $0.3$. The very end, near $k=30/7$, is also
certified. Two more spins are being computed now.
\end{itemize}

\textbf{How we checked it.}
\begin{itemize}\itemsep2pt
\item Every number in the certified parts was recomputed by a second, independent program.
\item Built-in ``must-fail'' tests showed the checks can say no.
\item Predictions were written down and committed before each long run, and every hit and miss is reported.
\item All code, outputs and logs are public.
\end{itemize}

\textbf{What we do \emph{not} claim.}
\begin{itemize}\itemsep2pt
\item That the whole range is a mathematical proof. Most of it is careful 50-digit numerics.
\item Anything about real astrophysical black holes. These spacetimes are not black holes, and nothing away from the
axis was checked.
\item Anything about other spins, other energies, or light rays (only four spins and one energy were tested).
\item That no conserved quantity of any kind exists. Only the well-behaved (``meromorphic'') kind is ruled out.
\item That an expert has checked this. The checks were done by separate AI sessions, not by human peer review.
\end{itemize}

\textbf{How it was made.} AI assistance (Claude by Anthropic for code, analysis and drafting; Gemini by Google for
critical reads) under the author's direction. The author takes full responsibility.

\textbf{Where to look.} Code and data: \href{https://doi.org/10.5281/zenodo.24593416}{doi:10.5281/zenodo.24593416};
\url{https://github.com/sumit7194/ansatz-machine}.
\end{document}
